% Options for packages loaded elsewhere
\PassOptionsToPackage{unicode}{hyperref}
\PassOptionsToPackage{hyphens}{url}
\documentclass[
  11pt,
]{article}
\usepackage{xcolor}
\usepackage[left=3cm,right=3cm,top=2.5cm,bottom=2.5cm]{geometry}
\usepackage{amsmath,amssymb}
\setcounter{secnumdepth}{5}
\usepackage{iftex}
\ifPDFTeX
  \usepackage[T1]{fontenc}
  \usepackage[utf8]{inputenc}
  \usepackage{textcomp} % provide euro and other symbols
\else % if luatex or xetex
  \usepackage{unicode-math} % this also loads fontspec
  \defaultfontfeatures{Scale=MatchLowercase}
  \defaultfontfeatures[\rmfamily]{Ligatures=TeX,Scale=1}
\fi
\usepackage{lmodern}
\ifPDFTeX\else
  % xetex/luatex font selection
\fi
% Use upquote if available, for straight quotes in verbatim environments
\IfFileExists{upquote.sty}{\usepackage{upquote}}{}
\IfFileExists{microtype.sty}{% use microtype if available
  \usepackage[]{microtype}
  \UseMicrotypeSet[protrusion]{basicmath} % disable protrusion for tt fonts
}{}
\usepackage{setspace}
\makeatletter
\@ifundefined{KOMAClassName}{% if non-KOMA class
  \IfFileExists{parskip.sty}{%
    \usepackage{parskip}
  }{% else
    \setlength{\parindent}{0pt}
    \setlength{\parskip}{6pt plus 2pt minus 1pt}}
}{% if KOMA class
  \KOMAoptions{parskip=half}}
\makeatother
\usepackage{longtable,booktabs,array}
\usepackage{caption}
\captionsetup[table]{skip=6pt}
\newcounter{none} % for unnumbered tables
\usepackage{calc} % for calculating minipage widths
% Correct order of tables after \paragraph or \subparagraph
\usepackage{etoolbox}
\makeatletter
\patchcmd\longtable{\par}{\if@noskipsec\mbox{}\fi\par}{}{}
\makeatother
% Allow footnotes in longtable head/foot
\IfFileExists{footnotehyper.sty}{\usepackage{footnotehyper}}{\usepackage{footnote}}
\makesavenoteenv{longtable}
\usepackage{graphicx}
\makeatletter
\newsavebox\pandoc@box
\newcommand*\pandocbounded[1]{% scales image to fit in text height/width
  \sbox\pandoc@box{#1}%
  \Gscale@div\@tempa{\textheight}{\dimexpr\ht\pandoc@box+\dp\pandoc@box\relax}%
  \Gscale@div\@tempb{\linewidth}{\wd\pandoc@box}%
  \ifdim\@tempb\p@<\@tempa\p@\let\@tempa\@tempb\fi% select the smaller of both
  \ifdim\@tempa\p@<\p@\scalebox{\@tempa}{\usebox\pandoc@box}%
  \else\usebox{\pandoc@box}%
  \fi%
}
% Set default figure placement to htbp
\def\fps@figure{htbp}
\makeatother
% definitions for citeproc citations
\NewDocumentCommand\citeproctext{}{}
\NewDocumentCommand\citeproc{mm}{%
  \begingroup\def\citeproctext{#2}\cite{#1}\endgroup}
\makeatletter
 % allow citations to break across lines
 \let\@cite@ofmt\@firstofone
 % avoid brackets around text for \cite:
 \def\@biblabel#1{}
 \def\@cite#1#2{{#1\if@tempswa , #2\fi}}
\makeatother
\newlength{\cslhangindent}
\setlength{\cslhangindent}{1.5em}
\newlength{\csllabelwidth}
\setlength{\csllabelwidth}{3em}
\newenvironment{CSLReferences}[2] % #1 hanging-indent, #2 entry-spacing
 {\begin{list}{}{%
  \setlength{\itemindent}{0pt}
  \setlength{\leftmargin}{0pt}
  \setlength{\parsep}{0pt}
  % turn on hanging indent if param 1 is 1
  \ifodd #1
   \setlength{\leftmargin}{\cslhangindent}
   \setlength{\itemindent}{-1\cslhangindent}
  \fi
  % set entry spacing
  \setlength{\itemsep}{#2\baselineskip}}}
 {\end{list}}
\usepackage{calc}
\newcommand{\CSLBlock}[1]{\hfill\break\parbox[t]{\linewidth}{\strut\ignorespaces#1\strut}}
\newcommand{\CSLLeftMargin}[1]{\parbox[t]{\csllabelwidth}{\strut#1\strut}}
\newcommand{\CSLRightInline}[1]{\parbox[t]{\linewidth - \csllabelwidth}{\strut#1\strut}}
\newcommand{\CSLIndent}[1]{\hspace{\cslhangindent}#1}
\setlength{\emergencystretch}{3em} % prevent overfull lines
\providecommand{\tightlist}{%
  \setlength{\itemsep}{0pt}\setlength{\parskip}{0pt}}
% Statistics in Medicine submission preamble.
% Shared across analysis/report/<paper>/report.Rmd files via
%   includes:
%     in_header: ../sim-preamble.tex
% in the YAML output block.
%
% The page footer (source path, render time, version) is no
% longer set here. It is supplied at render time by the vendored
% tools/stamp-render.R, which injects tools/stamp.tex.

% Continuous line numbering for editorial review. The 'mathlines'
% option extends numbering through display-math environments
% (equation, align, ...), which SIM editorial workflows expect.
\usepackage[mathlines]{lineno}
\linenumbers
\usepackage{amsmath}

% Spacing: linestretch is set in YAML (1.5 line spacing).
\usepackage{setspace}
%% tools/stamp.tex
%% Provenance footer for PDFs rendered from this project.
%% Vendored by zzcollab; this file is not project-specific. The
%% footer prints the source document and the time of rendering.
%% The git version is defined here too, and so is available to a
%% project that wants it on the page, but it is left off the
%% default footer: it already names the staged copy in share/ and
%% appears in its manifest row. All values are supplied at render
%% time by tools/stamp-render.R, which writes a small generated
%% file that redefines the macros below. The \providecommand
%% defaults keep a bare LaTeX compile working when the document is
%% built without the wrapper.
\usepackage{fancyhdr}
\usepackage{xcolor}
\providecommand{\stampsource}{\detokenize{(source not recorded)}}
\providecommand{\stamptime}{(time not recorded)}
\providecommand{\stampversion}{\detokenize{(version not recorded)}}
\setlength{\footskip}{40pt}
\newcommand{\stampfootcontent}{%
  \footnotesize\color{gray}%
  \parbox[b]{\textwidth}{%
    \stamptime\hfill\thepage\\[2pt]%
    \ttfamily\raggedright\stampsource}}
\pagestyle{fancy}
\fancyhf{}
\fancyfoot[C]{\stampfootcontent}
\renewcommand{\headrulewidth}{0pt}
\renewcommand{\footrulewidth}{0.4pt}
%% The title page uses the `plain` page style; redefine it so the
%% stamp appears there too.
\fancypagestyle{plain}{%
  \fancyhf{}%
  \fancyfoot[C]{\stampfootcontent}%
  \renewcommand{\headrulewidth}{0pt}%
  \renewcommand{\footrulewidth}{0.4pt}}
\renewcommand{\stampsource}{\textasciitilde{}/\allowbreak{}prj/\allowbreak{}res/\allowbreak{}36-pmsimstats-ng/\allowbreak{}pmsimstats-ng/\allowbreak{}analysis/\allowbreak{}report/\allowbreak{}01-dgp-mean-moderation-vs-mvn/\allowbreak{}report.Rmd}
\renewcommand{\stamptime}{2026-08-26 09:27 PDT}
\renewcommand{\stampversion}{\detokenize{39ae459-wip}}
\usepackage{amsmath}
\usepackage{amssymb}
\usepackage{booktabs}
\usepackage{longtable}
\usepackage{float}
\usepackage[mathlines]{lineno}
\linenumbers
\usepackage{bookmark}
\IfFileExists{xurl.sty}{\usepackage{xurl}}{} % add URL line breaks if available
\urlstyle{same}
% fallback for those not using the hyperref driver hyperxmp:
\makeatletter
\@ifundefined{xmpquote}{\newcommand{\xmpquote}[1]{#1}}{}
\makeatother
\hypersetup{
  pdftitle={Why carryover erodes power more under covariance moderation than mean moderation: a mechanistic analysis of biomarker-treatment interaction architectures in the Hybrid aggregated N-of-1 design},
  pdfauthor={pmsimstats team},
  hidelinks,
  pdfcreator={LaTeX via pandoc}}

\title{Why carryover erodes power more under covariance moderation than mean moderation: a mechanistic analysis of biomarker-treatment interaction architectures in the Hybrid aggregated N-of-1 design}
\author{pmsimstats team}
\date{August 26, 2026}

\begin{document}
\maketitle

{
\setcounter{tocdepth}{2}
\tableofcontents
}
\setstretch{1.1}
\section{Abstract}

\textbf{Background.} Simulation is the standard tool for power and
sample-size calculation in aggregated N-of-1 trials, designs in
which each participant is repeatedly crossed on and off treatment
and many single-patient series are pooled in order to validate a
predictive biomarker. Any such simulation must specify how the
biomarker moderates the treatment effect in its data-generating
process (DGP). That specification is a substantive choice rather
than an implementation detail, and each of the two established
parameterizations carries its own advantages and costs. We
formalize two contrasting parameterizations and ask why, in the
Hybrid N-of-1 design that motivates this software, carryover erodes
power so much more under one architecture than the other.

\textbf{Methods.} We define two DGP architectures. Under
`mean moderation' the biomarker scales the treatment effect in
the conditional mean. Under `covariance moderation' the
interaction is represented as a correlation that depends on
drug-exposure condition (on drug versus off drug), in a joint
multivariate normal (MVN) distribution. Holding the
analysis model fixed at a linear mixed model with continuous-time
AR(1) residual correlation, we quantify by Monte Carlo simulation
the power to detect the interaction as carryover increases in the
Hybrid N-of-1 design, the within-subject design used by the
prazosin and PTSD reference application that motivates this
software. A traditional crossover (CO) and an open-label-plus-
blinded-discontinuation design (OL+BDC) are run in parallel, at the
same total \(N\), as contrast cases that isolate which design feature
drives the architecture-dependent divergence.

\textbf{Results.} Under mean moderation power is nearly preserved
in the Hybrid design, because carryover does not disrupt the
biomarker's relationship to the treatment effect. At \(N = 70\),
power falls only from 0.842 to 0.784, a relative loss of 6.9\%
(\(z = 3.34\)), as the carryover half-life grows to one week. Under
covariance moderation the same carryover erodes the differential
correlation that carries the signal, and power falls from 0.840 to
0.711, a relative loss of 15.4\% (\(z = 2.80\), \(p = 0.005\)), more
than double the mean-moderation loss. The gap traces to a
two-channel erosion mechanism specific to covariance moderation:
carryover shrinks the drug-exposure contrast, a channel shared with
mean moderation, and independently decays the biomarker-response
correlation itself, a channel unique to covariance moderation. This
second channel's damage scales with how much uninterrupted off-drug
time a design exposes patients to, which is why the same
architecture contrast is negligible in the classical crossover (CO,
\(z = 0.29\), where within-subject AR(1) correlation already carries
most of the signal covariance moderation depends on) and far larger
in the open-label-plus-blinded-discontinuation design (OL+BDC,
\(z = 7.64\), \(p < 10^{-13}\), whose extended blinded-discontinuation
phase gives the correlation-decay channel the most uninterrupted
time to act). The Hybrid design's intermediate carryover exposure,
which alternates on- and off-drug periods rather than sustaining
either for long, places its architecture-dependent power loss
between these two extremes.

\textbf{Conclusions.} The DGP architecture is a substantive
scientific assumption rather than a parameterization convenience,
and it determines how severely carryover appears to reduce power.
Biomarker-validation simulation studies should state which of the
two architectures they adopt and report power under both.

\textbf{Keywords.} N-of-1 trials; biomarker-treatment interaction;
data-generating process; multivariate normal; carryover effects;
clinical trial simulation; precision medicine; statistical power.

\bigskip\noindent

\rule{\textwidth}{0.4pt}\bigskip

\section{1. Introduction}

Precision-medicine trials increasingly ask whether a baseline
biomarker modifies the treatment effect, that is, whether a
biomarker-treatment interaction is present. For simplified special
cases, such as a balanced strict repeated-measures analysis of
variance under sphericity, the power of the interaction test does
have a closed form, and a companion paper\textsuperscript{\citeproc{ref-pmsimstats-paper08}{1}}
derives it together with its equivalence to a two-sample \(t\)-test.
For the continuous-time mixed model and the N-of-1 designs
considered here, however, no such closed form is available, so
power is estimated by Monte Carlo simulation. Many trials are
generated and analyzed, and the rejection rate is recorded. The
simulation must itself specify the biomarker-treatment interaction
in its data-generating process, and there is more than one way to
do so.

We distinguish two architectures. The common one, direct
mean moderation, adds an interaction term to the outcome
equation so that the biomarker scales the size of the treatment
effect. Essentially every standard trial-simulation framework uses
it\textsuperscript{\citeproc{ref-simon2010}{2},\citeproc{ref-freidlin2014}{3}}. The second, covariance
moderation through differential correlation, used by Hendrickson
et al.\textsuperscript{\citeproc{ref-hendrickson2020}{4}} for N-of-1 trials, instead makes the
biomarker change the correlation between treatment response and
outcome, so that the interaction is represented in the joint
distribution rather than in the mean. On this reading it is the
reduced-form second-moment representation of a latent
responder-class mechanism, developed at length in a companion paper\textsuperscript{\citeproc{ref-pmsimstats-paper03}{5}}. Patients may belong to unobserved responder
and non-responder groups, with the biomarker only an imperfect
indicator of that status, so that two patients with the same
biomarker value can still respond differently. Mean moderation
cannot represent this situation, because it assumes that the
biomarker fixes each patient's drug effect exactly. Covariance
moderation represents it by letting biomarker and response track
one another only while the drug is acting, as a tendency rather
than a fixed rule, a tendency that fades once the drug washes out.
As Section 4 shows, that difference is why carryover costs more
power under covariance moderation than under mean moderation. A
third, `combined' architecture activates both channels
simultaneously through two independent parameters and recovers
these two as its single-channel limits. Because it introduces a
second free parameter and its power behavior warrants separate
treatment, it is developed in a companion paper\textsuperscript{\citeproc{ref-pmsimstats-paper11}{6}} and is not considered further here.

Choosing between these two architectures is not a matter of
computational convenience, because each carries genuine advantages
and a genuine cost. Mean moderation is simple, has a single free
parameter, and preserves the biomarker-treatment effect ratio under
carryover. It commits, however, to a strong assumption, namely that
the biomarker deterministically fixes each patient's drug response,
an assumption that is implausible whenever the biomarker is only a
statistical proxy for an unobserved responder subtype. Covariance
moderation relaxes that assumption. As the reduced-form
representation of a latent responder-class mechanism discussed
above, it lets the biomarker predict response only probabilistically,
which is the more defensible choice whenever responder status,
rather than the biomarker itself, is the true causal driver, as is
plausibly the case for the noradrenergic-tone hypothesis behind the
prazosin and PTSD reference application used throughout this paper.
That fidelity to the underlying biology comes at a real cost.
Because the covariance-moderation signal lives in a correlation
between biomarker and response that is sustained only while the
drug is active, any off-drug erosion of drug effect directly erodes
the statistical signal itself, producing power losses that this
paper shows can run many times larger than under mean moderation.
Investigators who adopt covariance moderation on biological grounds
should therefore budget for its carryover cost at the design stage,
rather than discover it after a trial has already been powered.

Hendrickson et al.\textsuperscript{\citeproc{ref-hendrickson2020}{4}} chose covariance moderation
for exactly this reason when they designed the Monte Carlo
simulation program behind the prazosin and PTSD trial. Given that
trial's own results, covariance moderation was the logical,
biologically appropriate specification, because representing the
interaction in the mean would have imposed a deterministic
biomarker-response relationship that the trial's biology did not
support. Mean moderation, by contrast, was and remains the modal
specification across the N-of-1 and broader biomarker-simulation
literature (Section 4.4), so Hendrickson's choice was a deliberate
departure from convention rather than an oversight. What that
departure costs, measured in statistical power under carryover, had
not previously been quantified, and neither had the mechanism
producing that cost been traced to specific, design-relevant
features of the trial. Doing both is the central aim of this paper.
We take the covariance-moderation choice as scientifically
justified for the prazosin and PTSD application, and we ask, in the
Hybrid design that trial and this software's target applications
use, why carryover erodes power so much more severely under
covariance moderation than under mean moderation, so that
investigators reasoning by analogy from that trial can budget for
the penalty explicitly and understand which design features control
its size.

In a single N-of-1 trial one participant is repeatedly crossed
between active treatment and a comparator over a sequence of
measurement occasions, so that the within-participant contrast
between on-drug and off-drug conditions estimates that
participant's individual treatment effect. An \emph{aggregated} N-of-1 trial pools many such
series in a mixed model to obtain population-level inference at far
smaller sample sizes than a parallel-group design. The feature that
distinguishes these designs analytically is \emph{carryover}:
residual drug effect that decays over days to weeks after
discontinuation and contaminates off-drug occasions. Our focus is
the Hybrid N-of-1 design, the within-subject design behind the
prazosin and PTSD reference application and the design this
software was built to support, which alternates on-drug and
off-drug periods rather than sustaining either for long. We run two
further within-subject designs in parallel as contrast cases that
isolate which feature of the Hybrid design's carryover exposure
drives its architecture-dependent power loss: a traditional
crossover (CO), whose within-subject AR(1) correlation dominates
any covariance-moderation signal, and an open-label-plus-blinded-
discontinuation design (OL+BDC), whose extended blinded phase gives
carryover the most uninterrupted time to act of the three. All
three designs are defined in Section 3.

\bigskip\noindent

\rule{\textwidth}{0.4pt}\bigskip

\section{2. Methods}

\subsection{2.1 Notation and analysis model}

Throughout, \(B_i\) denotes the biomarker value of participant \(i\),
\(D_{it}\) the binary on-drug indicator at timepoint \(t\) (1 on drug,
0 off drug), and \(BR_{it}\) the biological response component of the
outcome. These three symbols describe the data-generating-process
layer. The analysis-model layer uses a related but distinct
continuous drug-exposure variable, \texttt{Dbc}, the continuous
analogue of \(D_{it}\). It equals 1 during on-drug periods and
\(\exp(-\lambda t_{sd})\) during off-drug periods, where \(t_{sd}\) is
time since drug discontinuation and \(\lambda\) is the carryover
decay rate, typically derived from the drug's pharmacokinetic
half-life as \(\lambda = \ln 2 / t_{1/2}\). Whereas \(D_{it}\) flips
abruptly between 0 and 1, \texttt{Dbc} captures the smooth
shrinkage of residual drug effect during the off-drug period.
Code-style identifiers in typewriter font refer to variables in the
pmsimstats codebase or to R formula syntax. Mathematical symbols
are used otherwise. The analysis model is the same under both
architectures, an \texttt{nlme::lme} fit with a \texttt{corCAR1}
continuous-time autocorrelation structure, a random intercept, and
a single biomarker-by-drug interaction term, written
\texttt{bm:Dbc} in R formula syntax. The interaction coefficient on
\texttt{bm:Dbc} tests whether the biomarker moderates the
exposure-decayed drug effect.

\subsection{2.2 Two data-generating architectures}

\subsubsection{2.2.1 Mean moderation}

In this architecture the biomarker enters the mean structure of the
outcome directly. The treatment effect for participant \(i\) at
timepoint \(t\) is

\[Y_{it} = \mu_0 + \beta_1 \cdot D_{it} \cdot (1 + \beta_{bm}
\cdot B_i) + \epsilon_{it},\]

where \(D_{it}\) is the on-drug indicator (or a continuous measure of
drug exposure), \(B_i\) is the participant's biomarker value
(typically centered), and \(\beta_{bm}\) is the biomarker moderation
parameter. The interaction is explicit, deterministic, and lives in
the first moment of the outcome. Participants with higher biomarker
values have proportionally larger treatment effects, and this
population-level ratio is preserved under any proportional scaling
of \(D_{it}\), including the shrinkage that carryover imposes.

This multiplicative form maps directly onto the standard additive
regression form for predictive biomarker effects used in the trial
methodology literature, specifically the effect-modeling equation
of Kent et al.\textsuperscript{\citeproc{ref-kent2020path}{7}},
\(Y_{it} = \beta_0 + \beta_1 D_{it} + \beta_2 B_i + \beta_3 D_{it}
B_i + \epsilon_{it}\), with \(\beta_3 = \beta_1 \cdot \beta_{bm}\). The
additive form additionally accommodates a prognostic main effect
\(\beta_2\) for the biomarker, which the multiplicative form sets to
zero by construction. For power calculations on the interaction
coefficient the two forms are interchangeable up to this
prognostic-effect distinction, and both fall within mean
moderation.

\subsubsection{2.2.2 Covariance moderation through differential
correlation}

In this architecture the biomarker and the biological response (BR)
component are jointly drawn from a multivariate normal distribution
with a correlation that depends on drug-exposure condition,

\[\begin{pmatrix} B_i \\ BR_{it} \end{pmatrix} \sim
\text{MVN}\left(\begin{pmatrix} \mu_B \\ \mu_{BR}(t, D_{it})
\end{pmatrix}, \Sigma(D_{it})\right), \qquad
\text{Cor}(B_i, BR_{it}) = \begin{cases}
c_{bm} & D_{it} = 1 \\
c_{bm} \cdot e^{-\lambda t_{sd}} & D_{it} = 0,
\end{cases}\]

with the exponential-decay form as the default. The limiting cases
\(\lambda = 0\) (no decay) and \(\lambda \to \infty\) (immediate decay)
recover the two no-carryover model variants. The population mean of
\(BR_{it}\) is the same for all participants with the same treatment
history, so the interaction is not in the mean structure. It
instead emerges from the conditional distribution. Given \(B_i = b\),

\[E[BR_{it} \mid B_i = b, D_{it} = 1] = \mu_{BR}(t) + c_{bm} \cdot
\frac{\sigma_{BR}}{\sigma_B} \cdot (b - \mu_B).\]

The interaction here is probabilistic rather than deterministic,
lives in the second moment of the joint distribution, and persists
with attenuated strength off drug, vanishing only in the limit of
complete decay. Carryover in the DGP inflates the BR means during
off-drug periods, making them resemble on-drug means. This reduces
the differential in the correlation structure between drug-exposure
conditions. Section 2.2.3 shows that this erosion leaves the
analysis model's point estimate of the interaction largely
unbiased, but it converts variance the biomarker previously
explained into unexplained residual noise, which is the channel
through which it costs power (Section 3.2). A third, combined
architecture, in which both channels operate at
once through independent weights, recovers the two architectures
above as its single-channel limits and is the natural framework
when the operative mechanism is uncertain. It is developed in a
companion paper\textsuperscript{\citeproc{ref-pmsimstats-paper11}{6}}.

\subsubsection{2.2.3 Analytical comparison}

Consider the expected outcome difference between two participants
with biomarker values \(b_1\) and \(b_2\) at the same drug exposure
level \(D\). Under mean moderation,

\[E[Y \mid b_1, D] - E[Y \mid b_2, D] = \beta_1 \cdot \beta_{bm}
\cdot (b_1 - b_2) \cdot D.\]

This scales linearly with \(D\) and therefore shrinks during off-drug
periods, when \(D\) is replaced by an exposure-decayed value.
However, the \emph{ratio} of drug effect between high- and
low-biomarker participants is preserved at every level of exposure,
so the biomarker-treatment signal contracts in amplitude but does
not lose identifiability under carryover. Under covariance
moderation the analogous quantity, during off-drug periods with
carryover, is

\[E[BR \mid b, D=0, t_{sd}] = \mu_{BR,\text{off}}(t_{sd}) + c_{bm}
\cdot e^{-\lambda t_{sd}} \cdot \frac{\sigma_{BR}}{\sigma_B} \cdot
(b - \mu_B).\]

Here the biomarker-dependent component decays exponentially with
time off drug. As \(t_{sd}\) increases, the off-drug conditional
expectation converges to \(\mu_{BR,\text{off}}\) for all biomarker
values, in contrast to mean moderation's preserved ratio.

This convergence in the raw time since discontinuation is not, on
its own, the source of the power loss documented in Section 3, and
disentangling the two requires re-expressing the comparison in terms
of the regressor the analysis model actually uses. Write the
exponential decay factor as \(\kappa_{it} = 1\) on drug and
\(\kappa_{it} = e^{-\lambda t_{sd,it}}\) off drug, matching the
definition of \(D_{bc,it}\) in Section 2.1 exactly when the
analysis-model half-life is set equal to the DGP half-life, the base
case examined in Section 3.1. Then \(\text{Cor}(B_i, BR_{it}) =
c_{bm} \cdot D_{bc,it}\) by construction, and the conditional mean
above is \emph{exactly linear} in \(D_{bc,it} \cdot (b - \mu_B)\) at
every value of \(t_{sd}\), not only in the untouched limit:
\[
E[BR_{it} \mid B_i = b, D_{it}] = \mu_{BR}(t) + c_{bm}\, D_{bc,it}\,
\frac{\sigma_{BR}}{\sigma_B} (b - \mu_B).
\]
Because the analysis model regresses the outcome on
\texttt{bm:Dbc} using this same \(D_{bc,it}\), its population
(large-sample) estimate of the interaction slope recovers
\(c_{bm}\,\sigma_{BR}/\sigma_B\) \emph{regardless of how much
carryover has occurred}, to the same leading order that mean
moderation's slope \(\beta_1\beta_{bm}\) is carryover-invariant. The
biomarker-dependent signal is not eliminated by carryover in the
sense that matters for the fitted interaction coefficient; only its
raw amplitude at long \(t_{sd}\) is. This is confirmed empirically in
Section 3.4: the mean of \(\hat\beta_{bm:D_{bc}}\) under covariance
moderation holds near \(-0.234\) across every carryover level examined,
rather than trending toward zero, exactly as this identity predicts.

If the population slope is unaffected, the power loss under
covariance moderation must arise instead from the precision with
which that slope is estimated, and the bivariate normal structure of
Section 2.2.2 gives that precision loss in closed form. Conditional
on \(B_i\), the variance of \(BR_{it}\) is
\[
\mathrm{Var}(BR_{it} \mid B_i, D_{it}) = \sigma_{BR}^2
\bigl(1 - c_{bm}^2\, D_{bc,it}^2\bigr),
\]
an exact consequence of the bivariate normal conditional-variance
identity applied to a correlation of \(c_{bm}\, D_{bc,it}\). On drug
(\(D_{bc,it} = 1\)), the biomarker explains a fixed \(c_{bm}^2\) fraction
of \(BR\)'s variance; as carryover erodes \(D_{bc,it}\) toward zero off
drug, that explained fraction shrinks \emph{quadratically}, not
linearly, in \(D_{bc,it}\), and the corresponding conditional variance
rises toward the full unconditional \(\sigma_{BR}^2\), adding pure
noise that the fixed-effects analysis model, which assumes
homoscedastic residual variance, cannot distinguish from
measurement error. Mean moderation has no analogue of this channel:
its residual variance \(\epsilon_{it}\) is constant by construction
(Section 2.2.1), independent of \(D_{it}\), so carryover there degrades
only the regressor's contrast, never the noise floor around it. This
is the quantitative content behind the qualitative two-channel
account of Section 3.2: mean moderation loses information only
through the design contrast in \(D_{bc,it}\) itself, a channel shared
by both architectures, while covariance moderation loses information
through this quadratic variance-inflation channel as well, and it is
this second, architecture-specific channel that produces the
divergence in power documented in Section 3.

\subsection{2.3 Simulation design}

We ran identical simulation configurations under both architectures
using the pmsimstats-ng framework, matching the DGP parameters as
closely as possible. The non-biologic response components
(time-varying response, placebo-belief accumulation, and residual
noise) are generated from the same multivariate normal draw, with
identical means, variances, within-factor AR(1) autocorrelation,
and cross-factor correlations, under both architectures. The two
DGPs differ only in how the biomarker-to-biological-response
linkage is specified, so the divergent carryover behavior reported
in Section 3 is attributable to that parameterization alone and not
to any difference in how the other components are generated.

We follow the ADEMP framework of Morris, White, and Crowther\textsuperscript{\citeproc{ref-morris2019simulation}{8}}, comprising Aims, Data-generating mechanisms,
Estimands, Methods, and Performance measures. The Aim is to compare
power to detect the biomarker-by-drug interaction across
architectures as carryover increases. The Data-generating
mechanisms are the two architectures of Section 2.2. The
\textbf{primary estimand} is the interaction coefficient
\(\beta_{bm:D}\) in the model
\texttt{Sx \textasciitilde{} bm + t + Dbc + bm:Dbc}. The
\textbf{Methods} setup is as follows.

\begin{itemize}\setlength{\itemsep}{2pt}\setlength{\parskip}{0pt}
\item Trial designs: CO (traditional crossover), Hybrid (N-of-1
  Hybrid), OL+BDC (open-label with blinded discontinuation).
\item Sample size: $N = 70$ in total for every design, allocated
  across each design's randomization paths, so the three designs
  are matched on $N$. The two-path CO and OL+BDC designs put 35 on
  each path, and the four-path Hybrid design splits 70 as
  18/18/17/17.
\item Biomarker correlation or moderation: $c_{bm} = 0.45$,
  equivalently $\beta_{bm} = 0.45$.
\item DGP carryover half-life: $t_{1/2} \in \{0, 0.5, 1.0\}$ weeks,
  with the analysis-model \texttt{Dbc} half-life matched to it.
\item Replicates: 1{,}000 per cell (Monte Carlo standard error, MCSE,
  on a binomial power estimate near 0.5 is approximately 0.016).
\item Implementation: 8-core \texttt{parallel::mclapply}, full
  factorial of 2 architectures $\times$ 3 designs $\times$ 3
  $t_{1/2}$ $\times$ 2 $c_{bm}$ levels gives 36 cells and 36{,}000
  fits, with 100\% convergence.
\end{itemize}

The \textbf{Performance measures} are power for \(\beta_{bm:D}\) at
\(\alpha = 0.05\), Type I error under \(c_{bm} = 0\) (both with MCSE
\(\sqrt{\hat\pi(1-\hat\pi)/n_\text{reps}}\)), bias of
\(\hat\beta_{bm:D}\) relative to the reference true value
\(\theta_\text{true} = -c_{bm} \cdot \sigma_{BR}\), the same
calibration used in the companion carryover-specification and
component-decomposition manuscripts\textsuperscript{\citeproc{ref-pmsimstats-paper02}{9},\citeproc{ref-pmsimstats-paper06}{10}}, the empirical MCSE of \(\hat\beta_{bm:D}\) across
replicates, and the convergence rate. The number of replicates per
cell is chosen to achieve MCSE below 0.02 for power at \(\hat\pi =
0.5\), which corresponds to \(n_\text{reps} \geq 625\). We use
\(n_\text{reps} = 1000\) throughout.

The \(N = 70\) choice is the smallest of the prazosin-PTSD-matched
sample sizes used in the published methodological literature, and
it is a regime in which no design or architecture saturates power
at the no-carryover cells (the highest cell in the grid is 0.842),
so carryover-induced erosion is mechanically visible. This is the
main practical gain from matching the designs on \(N\). At a previous
per-path parameterization the Hybrid cells sat above 0.99, where no
erosion could register.

\bigskip\noindent

\rule{\textwidth}{0.4pt}\bigskip

\section{3. Results}

Covariance moderation, as Section 1 argues, is often the more
biologically defensible choice when a biomarker is best understood
as predictive rather than causal, and it is the specification
Hendrickson et al.~selected for the prazosin and PTSD application on
exactly these grounds. That defensibility does not eliminate a
carryover cost. Quantifying that cost in the Hybrid design, and
tracing why it arises, so that the appeal of the architecture does
not obscure a real design consideration, is the purpose of what
follows. The crossover (CO) and open-label-plus-blinded-
discontinuation (OL+BDC) designs are reported alongside Hybrid at
every step, not as objects of primary interest, but because
contrasting Hybrid against these two structurally different designs
is what identifies which feature of a design's carryover exposure
drives the architecture-dependent loss.

\subsection{3.1 Power under carryover at \emph{N} = 70}

\textbf{Mean moderation.}

\begin{table}[htbp]
\centering
\begin{tabular}{lccc}
\toprule
Design & $t_{1/2} = 0$ & $t_{1/2} = 0.5$ & $t_{1/2} = 1.0$ \\
\midrule
CO & 0.739 & 0.739 & 0.725 \\
Hybrid & 0.842 & 0.844 & 0.784 \\
OL+BDC & 0.751 & 0.748 & 0.730 \\
\bottomrule
\end{tabular}
\end{table}

Power declines modestly under mean moderation. In the Hybrid
design, the relative loss from \(t_{1/2} = 0\) to \(t_{1/2} = 1\) is
6.9\%, and it is the only one of the three designs whose loss is
separable from zero at 1\{,\}000 replicates per cell (\(z = 3.34\)).
The two-period designs lose less (1.9\%, CO; 2.8\%, OL+BDC) and
neither loss is distinguishable from Monte Carlo noise (CO
\(z = 0.71\); OL+BDC \(z = 1.07\)). All three values fall at or below
the 8-13\% relative-loss range suggested by earlier exploratory
runs in this project for direct-mean-moderation DGPs, an erosion
that leaves power substantially intact even in the design, Hybrid,
where it is largest.

\textbf{Covariance moderation (differential correlation, MVN).}

\begin{table}[htbp]
\centering
\begin{tabular}{lccc}
\toprule
Design & $t_{1/2} = 0$ & $t_{1/2} = 0.5$ & $t_{1/2} = 1.0$ \\
\midrule
CO & 0.745 & 0.754 & 0.723 \\
Hybrid & 0.840 & 0.795 & 0.711 \\
OL+BDC & 0.777 & 0.694 & 0.539 \\
\bottomrule
\end{tabular}
\end{table}

Power declines more sharply under covariance moderation than under
mean moderation in every design, but by very different amounts. In
the Hybrid design, the relative loss from \(t_{1/2} = 0\) to
\(t_{1/2} = 1\) is 15.4\%, more than double its matched
mean-moderation loss (6.9\%). The two-period designs bracket this
value from opposite sides: CO loses only 3.0\% (not significantly
different from mean moderation's CO loss), while OL+BDC loses
\textbf{30.6\%}, roughly 11 times its matched mean-moderation loss
(2.8\%) and twice the Hybrid loss. The CO cell shows a small,
non-monotone fluctuation (\(0.745 \to 0.754 \to 0.723\)), within one
MCSE and consistent with no architecture-by-carryover interaction in
that design. These figures fall below the 40-60\% relative-loss
range suggested by earlier exploratory runs for the MVN
differential-correlation DGP, but confirm their qualitative finding
of substantial, design-dependent erosion, with the Hybrid design
occupying a genuine middle position rather than resembling either
contrast case.

We compare the carryover-induced power loss between architectures
using a difference-of-differences \(z\)-statistic. Let
\(\Delta_X = \hat\pi_{X,0} - \hat\pi_{X,1}\) denote the
within-architecture loss from \(t_{1/2} = 0\) to \(t_{1/2} = 1.0\) for
architecture \(X\), estimated from \(n = 1{,}000\) independent
replicates per cell. The covariance-minus-mean gap
\(\Delta_B - \Delta_A\) has standard error

\[
\text{SE} = \sqrt{
  \frac{\hat\pi_{B,0}(1-\hat\pi_{B,0})}{n}
+ \frac{\hat\pi_{B,1}(1-\hat\pi_{B,1})}{n}
+ \frac{\hat\pi_{A,0}(1-\hat\pi_{A,0})}{n}
+ \frac{\hat\pi_{A,1}(1-\hat\pi_{A,1})}{n}},
\]

since the four proportions come from independent replicate sets
(all tests two-sided). In the Hybrid design, \(\Delta_A = 0.842 -
0.784 = 0.058\) and \(\Delta_B = 0.840 - 0.711 = 0.129\), giving
\(\Delta_B - \Delta_A = 0.071\) (about 7 percentage points) and
\(z = 2.80\) (\(p = 0.005\)): the architectures' carryover behavior
differs well beyond Monte Carlo error even in this intermediate
design. The same contrast is much larger for OL+BDC
(\(\Delta_A = 0.021\), \(\Delta_B = 0.238\), \(\Delta_B - \Delta_A =
0.217\), about 22 percentage points, \(z = 7.64\), \(p < 10^{-13}\)) and
indistinguishable from zero for CO (\(z = 0.29\)), because the CO
design's within-subject AR(1) structure dominates the
carryover-mediated correlation-decay channel that covariance
moderation uses to represent the interaction, leaving little
architecture-dependent signal for the contrast to detect. The
ordering of the three designs on this contrast, OL+BDC well above
Hybrid and Hybrid above CO, is the substantive result, and it is
consistent with a single explanatory factor: the amount of
uninterrupted off-drug time each design's structure exposes
patients to. It is read at a common \(N = 70\) and is not a statement
about the designs' relative power (Section 4.1).

\subsection{3.2 Mechanism: why the architectures diverge}

Section 2.2.3 identified two channels through which carryover
degrades the interaction signal, and only one of them is shared by
both architectures. \textbf{Design-contrast erosion} shrinks
\(\mathrm{Var}(D_{bc,it})\) as carryover pulls off-drug exposure values
away from a clean 0/1 contrast; because the interaction regressor is
\texttt{bm:Dbc}, less contrast in \texttt{Dbc} means less
information for the interaction coefficient under either
architecture, and this channel alone explains mean moderation's
modest, design-dependent loss (1.9-6.9\%). \textbf{Conditional-variance
inflation} is unique to covariance moderation and is the channel
that produces the much larger losses of Section 3.1. As
\(D_{bc,it} \to 0\) off drug, \(\mathrm{Var}(BR_{it} \mid B_i, D_{it}) =
\sigma_{BR}^2(1 - c_{bm}^2 D_{bc,it}^2)\) rises from its on-drug floor
of \(\sigma_{BR}^2(1-c_{bm}^2)\) toward the full unconditional
\(\sigma_{BR}^2\), converting variance that was explained by the
biomarker into unexplained residual noise that a homoscedastic
fixed-effects model cannot separate from measurement error. Mean
moderation has no such channel: its residual variance is constant by
construction, so it loses only design contrast, never precision
around a fixed contrast. This is why the population interaction
coefficient stays essentially unbiased under both architectures
(Section 3.3) while power diverges sharply: the loss is a precision
effect specific to covariance moderation, not a bias effect shared by
both.

A rough, working-independence approximation to the interaction
test's noncentrality parameter makes the connection to design and
carryover explicit. Treating each observation's contribution to the
Wald statistic as approximately additive and ignoring the
\texttt{corCAR1} residual correlation that the actual analysis model
estimates, the information available from an observation at exposure
level \(D_{bc,it} = x\), after partialling out the model's own
\texttt{Dbc} main effect, is approximately proportional to
\((x - \bar D_{bc})^2/\sigma_\epsilon^2\) under mean moderation and to
\((x - \bar D_{bc})^2 / \bigl[\sigma_{BR}^2(1 - c_{bm}^2
x^2)\bigr]\) under covariance moderation, where \(\bar D_{bc}\) is the
design's mean exposure level. Both summands shrink as carryover
compresses \(D_{bc,it}\) toward \(\bar D_{bc}\) across the design, and
the covariance-moderation summand is additionally discounted by the
same conditional-variance-inflation factor derived above, so the
total information lost accumulates fastest in designs that spend the
most observation-weeks away from a clean on/off contrast.

This approximation was evaluated numerically against each of the
three designs' exact timepoint and phase schedules (constructed with
the package's own \texttt{buildtrialdesign()}, at \(c_{bm} = 0.45\)),
rather than left as an unexecuted next step, and the result is mixed.
Summing \((D_{bc,it} - \bar D_{bc})^2\) over every observation in each
design correctly orders the three by predicted information loss from
\(t_{1/2} = 0\) to \(t_{1/2} = 1.0\): 36.8\% in the focal Hybrid design,
against 4.6\% in CO and 51.2\% in OL+BDC, matching the qualitative
ordering of Section 3.1's power losses under either architecture.
But the covariance-moderation summand predicts an architecture-specific
gap over mean moderation's own loss of under one percentage point at
every design (37.1\% in Hybrid, 4.7\% in CO, 50.9\% in OL+BDC) at
\(c_{bm} = 0.45\), far smaller than the large gap Section 3.1 actually
reports (for example, Hybrid's 6.9\% mean-moderation loss against its
15.4\% covariance-moderation loss, or OL+BDC's 2.8\% against 30.6\%).
The exact conditional-variance-inflation identity itself is not in
question here; it is independently confirmed against realized Monte
Carlo variability in Section 3.3. What this numeric check shows is
that the working-independence simplification used to connect that
identity to design and carryover analytically discards enough of the
\texttt{corCAR1} structure to badly underpredict the size of the
architecture gap, not only its near-absence in CO. A full analytical
account of the magnitude, rather than only its direction, would need
to incorporate the within-subject autocorrelation directly, which
remains undone; the qualitative account below is offered as the
mechanism for CO's particular robustness in place of that missing
derivation.

Qualitatively, the ordering of Section 3.1 follows the same
information argument applied to each design's own pattern of
drug-exposure alternation. Conditional-variance inflation needs
sustained low-\(D_{bc,it}\) time to accumulate: the longer a design
lets \(D_{bc,it}\) sit near zero before the next on-drug period resets
it, the more of that period's information is converted to noise.
OL+BDC sustains its off-drug window the longest of the three
designs, so its variance-inflation channel does the most damage. CO
alternates quickly enough, and leans on within-subject AR(1)
correlation heavily enough to have already extracted most of the
available signal by other means, that little conditional-variance
inflation has room to act before the next on-drug period arrives.
The Hybrid design's paths mix on- and off-drug periods at an
intermediate cadence, so \(D_{bc,it}\) rarely stays near zero as long
as it does under OL+BDC, capping how much variance inflation any
single off-drug stretch can accumulate; the design's intermediate
architecture-dependent loss (Section 3.1) is the empirical signature
of this intermediate exposure pattern.

\subsection{3.3 Type I error and estimator bias}

Across all 18 null cells of the Section 3.1 production run, Type I
error sat in \([0.010, 0.074]\), with mean 0.042 and median 0.035, and
no architecture-specific inflation. The uniformly sub-nominal rates
are consistent with the \texttt{corCAR1} working-covariance
miscalibration characterized in a companion calibration study\textsuperscript{\citeproc{ref-pmsimstats-paper10}{11}}, whose model-based standard error is
over-stated by roughly 6 to 10 percent. Because this affects both
architectures alike, it is common-mode with respect to the
architecture contrast and does not affect it.

Mean estimated \(\hat\beta_{bm:D_{bc}}\) at \(c_{bm} = 0.45\) ranged
from \(-0.230\) to \(-0.281\) across the 18 alternative cells, with the
mean-moderation OL+BDC cell at \(t_{1/2} = 1.0\) producing the
largest absolute mean. Covariance-moderation cells maintained a
mean near \(-0.234\) across all \(t_{1/2}\) values, exactly as the
unbiasedness argument of Section 2.2.3 predicts: the population
interaction slope recovered from \(D_{bc,it}\) does not depend on
carryover once the analysis model's half-life matches the DGP's.
Empirical SE of \(\hat\beta_{bm:D_{bc}}\) runs 0.07 to 0.10 across
cells, and the within-replicate model SE matches it at scale, which
is consistent with nominal-95\% Wald coverage but was not directly
verified by a coverage simulation.

The conditional-variance-inflation mechanism of Section 3.2 can be
checked directly against this same table, since it predicts where
the empirical SE should grow with carryover and where it should not.
Computing \(|\hat\beta_{bm:D_{bc}}|/\mathrm{SE}(\hat\beta_{bm:D_{bc}})\)
at \(c_{bm} = 0.45\) as a design- and architecture-specific proxy for
the interaction test's signal-to-noise ratio, and comparing its value
at \(t_{1/2} = 0\) against \(t_{1/2} = 1.0\): under covariance moderation
this ratio falls by 12.2\% in OL+BDC (2.947 to 2.588) and 12.6\% in
Hybrid (3.259 to 2.849), while under mean moderation it falls by only
5.4\% in Hybrid (3.002 to 2.839) and is flat to slightly higher in
OL+BDC (\(-2.7\%\), 2.865 to 2.943) and in CO (0.6\%, mvn; 1.9\%, mean
moderation). This is the pattern Section 3.2 predicts: covariance
moderation's SNR erosion is driven almost entirely by SE growth
(0.079 to 0.090 in OL+BDC, a 14.6\% increase) against an essentially
flat numerator, consistent with the derived
\(\sigma_{BR}^2(1 - c_{bm}^2 D_{bc,it}^2)\) variance-inflation
identity, whereas mean moderation shows no comparable SE trend
because its residual variance does not depend on \(D_{bc,it}\). This
check is verified directly from the committed cell-level summaries
(\texttt{01-dgp-summary.txt}) rather than from a new simulation, and
it substantiates the mechanism of Section 3.2 quantitatively, though
it remains a post hoc check on realized Monte Carlo variability
rather than a derivation of the noncentrality parameter from design
and DGP parameters alone.

\bigskip\noindent

\rule{\textwidth}{0.4pt}\bigskip

\section{4. Discussion}

\subsection{4.1 Matching architecture to biological mechanism}

The results of Section 3 reduce to a single trade-off that this
subsection develops at length. Mean moderation is simpler, is
essentially immune to carryover, and is directly comparable to the
overwhelming majority of the published biomarker-simulation
literature, but it assumes a mechanism that is often biologically
wrong: that the biomarker fixes each patient's drug response
exactly, rather than merely predicting it. Covariance moderation is
biologically faithful in exactly the cases where mean moderation is
not, because it lets the biomarker predict response only
probabilistically, but that fidelity is purchased at a real,
design-dependent power cost under carryover, a cost of the kind that
an investigator who ran their power calculation under mean
moderation alone would never see in their own simulation and would
therefore not budget for. Neither architecture dominates the other
in the abstract; the correct choice, and the size of the risk from
choosing wrong, depends on the biology of the specific biomarker
under study, a claim this subsection substantiates first from first
principles and then from the finite-mixture formalism developed at
length in a companion paper.

The choice of architecture is a choice about \emph{why} the
biomarker predicts response, and is therefore a biological
statement rather than a coding preference. Mean moderation is
appropriate when the biomarker mechanically sets the size of the
drug effect, for example by fixing effective dose or receptor
occupancy, so that doubling the biomarker roughly doubles the
effect at every timepoint. Because carryover keeps the same
proportion off drug, power is barely touched. Typical examples
include drug clearance (kidney function setting the steady-state
level of a renally cleared drug), receptor density, and metabolizer
genotype. Covariance moderation is appropriate when the biomarker
is only a proxy for a hidden responder status that it indexes
imperfectly, so that identical biomarker values can still respond
differently, and the biomarker predicts response only while the
drug is acting. Once the drug washes out, the correlation that
carried the signal fades; the analysis model's point estimate stays
close to unbiased even so (Section 2.2.3), but the precision with
which it can be estimated drops, and power falls sharply as a
result (Section 3.2). Typical examples include resting blood
pressure as a marker of
noradrenergic tone, inflammatory markers such as CRP or IL-6, and
polygenic risk scores. A fully faithful model of the hidden
responder-class idea would be a finite mixture. Covariance
moderation is a convenient correlation-level approximation, adequate
for power calculations but not a claim about the true mechanism. A
companion paper\textsuperscript{\citeproc{ref-pmsimstats-paper03}{5}} develops the mixture and
latent-class formulations, deriving when a single multivariate
normal reproduces a two-component mixture's first two moments and
cataloguing where it fails. Senn\textsuperscript{\citeproc{ref-senn2016mastering}{12}} also cautions
that apparent biomarker-treatment effects are sometimes just
variance artifacts, so covariance moderation's signal may be
smaller than assumed.

The pmsimstats software was built for a prazosin/PTSD trial\textsuperscript{\citeproc{ref-hendrickson2020}{4},\citeproc{ref-raskind2013}{13}}, in which the biomarker, resting
blood pressure, is thought to be a proxy for noradrenergic tone
rather than a direct driver of drug levels, a covariance-moderation
story. The carryover-driven power loss we find under covariance
moderation is therefore a clinically relevant warning. Designs with
long off-drug periods may lose more power to detect this biomarker
than a mean-moderation simulation would suggest. Should blood
pressure act through the mean channel as well as the correlation
channel, the dual-channel architecture of the companion paper\textsuperscript{\citeproc{ref-pmsimstats-paper11}{6}} is the appropriate model, and because it
estimates two independent parameters, it calls for independent
evidence on each.

Three further considerations support covariance moderation as the
more defensible specification for this application. First, resting
blood pressure is a single-occasion vital sign subject to
substantial physiological and measurement variability, so treating
it as a precise multiplicative scale factor on the treatment effect
assumes a level of measurement precision the biomarker does not
have. A specification that asserts only an association, and its
strength, is proportionate to that measurement error. Second, the
distinction tracks an established one in the predictive-biomarker
literature. The PATH framework of Kent et al.\textsuperscript{\citeproc{ref-kent2020path}{7}}
separates a causal treatment effect modifier from a biomarker that
is merely predictive of response without being causally upstream of
it. Covariance moderation is the natural representation of the
predictive role and mean moderation of the causal-mediator role, so
classifying blood pressure as predictive rather than causal in this
trial argues for covariance moderation on grounds already
established in that literature. Third, mean moderation requires
committing to a specific functional form for how the biomarker
scales the treatment effect, typically linear in the centered
biomarker. Covariance moderation makes the weaker assumption that
biomarker and response are associated, without specifying the
functional form of that association beyond its strength, which is
the more parsimonious choice when investigators lack direct
evidence for a particular dose-response shape.

The two architectures considered here are the two poles of a wider
spectrum, and drug and biomarker properties can place a candidate
application closer to one boundary of that spectrum than to either
pole outright. A companion paper on latent-class and mixture-model
formulations\textsuperscript{\citeproc{ref-pmsimstats-paper03}{5}} develops this spectrum at length
and grounds covariance moderation as a specific, derivable
approximation to a genuinely latent-class generative mechanism, an
argument summarized here because it is the strongest available
justification for covariance moderation's biological
faithfulness and, at the same time, the source of its principal
limitation.

The mixture argument begins from the observation that if a
candidate biomarker indexes an unobserved responder or
non-responder dichotomy, as the noradrenergic-tone hypothesis
behind the prazosin and PTSD reference application supposes, the
mechanistically faithful DGP is not a correlation but a finite
mixture,
\[
Z_i \mid B_i \sim \mathrm{Bernoulli}\bigl(\pi(B_i)\bigr), \qquad
BR_{it} \mid Z_i = z \sim f_z(t, D_{bc,it}), \qquad
\pi(B_i) = \mathrm{expit}(\psi_0 + \psi_g\, b_i),
\]
where \(Z_i\) is an unobserved class label fixed for participant \(i\)
across the whole trial, \(b_i\) is the standardized biomarker, and the
gating slope \(\psi_g\) governs how steeply the biomarker predicts
class membership. Two properties of this generative form matter for
what follows. First, \(Z_i\) is invariant across drug state and time:
once a participant is drawn into class 1 or class 0, that
membership persists through every on-drug and off-drug measurement
occasion in the trial, in principle recoverable from the full
within-subject trajectory even at timepoints where the
instantaneous drug-exposure contrast has decayed to near zero.
Second, the class-conditional mean trajectory itself still depends
on \(D_{bc,it}\), so a full carryover off-drug period drives the two
class-conditional means toward each other if the responder and
non-responder distinction lies purely in drug response, exactly the
mean-blurring channel that Section 3.2 identifies as common to both
architectures examined here.

Under within-class independence between \(B_i\) and \(BR_{it}\), a
two-component mixture of this form induces a marginal joint
distribution of \((B_i, BR_{it})\) whose first two moments are
reproduced, to leading order, by a single multivariate normal with
\(\mathrm{Cor}(B, BR) = c_{bm}\), where
\[
c_{bm}^2 = \frac{\pi(1-\pi)(\mu_B^1 - \mu_B^0)^2}{\mathrm{Var}(B)}
\cdot
\frac{\pi(1-\pi)(\mu_{BR}^1 - \mu_{BR}^0)^2}{\mathrm{Var}(BR)},
\]
the product of the between-class variance fractions in the
biomarker and in the biological-response component. Under mild
regularity conditions (adequate class overlap, mixing proportion
\(\pi\) bounded away from 0 and 1), this second-moment structure
carries a differential correlation across drug-exposure conditions,
elevated where the class-conditional means separate on drug and
shrinking toward zero where they coincide off drug after carryover
has run its course, which is precisely the covariance-moderation DGP
of Section 2.2.2. Covariance moderation is therefore not an
arbitrary alternative to mean moderation; it is the specific
second-moment approximation to a latent-class mechanism that a
biomarker acting as an imperfect proxy for a hidden responder
subtype would actually generate, which is the formal grounding for
treating it as the more biologically faithful choice whenever that
subtype story is plausible.

That derivation also exposes covariance moderation's principal
limitation directly. An MVN approximation reproduces only the first
two moments of the mixture; it discards the class label \(Z_i\)
itself, along with the class-invariance property that made \(Z_i\)
carryover-robust in the true generative form. A correlation, unlike
a class label, must be estimated fresh from the currently observed
joint distribution at every timepoint, with no persistent
per-participant memory once off-drug attenuation has taken effect.
This is why the power losses documented in Section 3 are a property
of the covariance-moderation \emph{approximation} specifically, not
an inherent, unrecoverable cost of a genuinely latent-class biology.
A class-aware analysis fitted to mixture-DGP data could, in
principle, recover some of that apparent loss by using the
carryover-invariant class label directly, up to an asymptotic bound
set by the gap between the covariance-only and mixture-faithful
power estimates. In practice, though, that recovery is bounded by
identifiability: fitting a finite mixture (an EM algorithm, or its
implementations in \texttt{lcmm} or \texttt{flexmix}) reliably
distinguishes genuine class structure from noise only at sample
sizes of several hundred to several thousand in the psychometric
literature that developed these methods\textsuperscript{\citeproc{ref-bauercurran2003}{14}--\citeproc{ref-nylund2007}{16}}, well above the
\(N \in [30, 150]\) typical of N-of-1 biomarker-validation trials. At
trial-relevant \(N\), mixture analysis may be underpowered or
unidentified even when the underlying biology is genuinely
mixture-structured, so a class-aware primary analysis is not
generally a usable alternative to the standard mixed-model analysis
of Section 2.1, only a sensitivity check where identifiability
permits. Covariance moderation is therefore best understood as the
pragmatic compromise this identifiability constraint forces: the
representation that is faithful to the underlying biology at the
level of the joint distribution's first two moments, without
requiring an estimation strategy that trial-relevant sample sizes
cannot support. When the biomarker is a
near-deterministic class indicator, for example a genotype with
near-perfect sensitivity and specificity for a metabolizing versus
non-metabolizing phenotype, the class-membership gradient is steep
and a thresholded mean-moderation specification is the better
approximation even though the underlying mechanism is a latent
subtype rather than a graded scaling. When the two response classes
are well separated, as in some targeted-oncology indications with a
driver mutation that essentially abolishes response in its absence,
the marginal response distribution is genuinely bimodal and neither
architecture is adequate: a full mixture DGP, not examined by
simulation in this paper, is required to avoid understating the
power achievable by a class-aware analysis. When covariance,
autocorrelation, or placebo-response parameters differ by latent
class rather than only the mean or the biomarker correlation, a
single-component MVN of either architecture is misspecified
regardless of which moment carries the interaction. Investigators
should therefore treat mean moderation and covariance moderation as
convenient limiting cases to be checked against, rather than an
exhaustive pair of options, whenever pilot data suggest the
biomarker may index a genuinely discrete or highly separated
subgroup structure.

The trial-design dependence documented in Section 3, largest in
OL+BDC and smallest in CO, is the empirical signature of the same
class-label discarding just established: because covariance
moderation carries no persistent per-participant memory of
responder status once off-drug attenuation has run its course, the
signal it can offer a standard analysis shrinks with every
additional unit of uninterrupted off-drug time a design imposes, in
exactly the pattern Section 3.2 traces mechanistically across the
three designs. So the design implication drawn in Section 4.4, that
designs with long off-drug windows carry the largest
architecture-dependent risk, should be read as a caution for trials
analyzed with the standard mixed-model methods of Section 2.1
specifically, not as a statement that latent-class biology makes
long off-drug windows safe by construction; a class-aware analysis,
where trial-relevant identifiability permits one, is not bound by
the same limitation. For the same reason, the ordering of the three
designs on architecture-sensitivity (Section 3.1) is a statement
about how much power a fixed analysis model loses to the architecture
choice, not a design-preference recommendation. At matched \(N = 70\),
covariance-moderation power under carryover happens to land close
together for Hybrid (0.711) and CO (0.723) and well below both for
OL+BDC (0.539), but the paths are allocated unevenly across the
designs (18/18/17/17 for Hybrid's four paths versus 35/35 for CO's
and OL+BDC's two), and design selection in practice turns on
feasibility, blinding, and regulatory considerations that this power
comparison alone does not capture. Investigators choosing among
designs on carryover-robustness grounds should treat the Section
3.1 ordering as a hypothesis to examine under their own design's
path allocation, not as a ranking to adopt directly.

\subsection{4.2 Reporting recommendations}

\begin{table}[htbp]
\centering
\begin{tabular}{lcc}
\toprule
Criterion & Mean moderation & Covariance moderation \\
\midrule
Biomarker role & Causal mediator & Statistical predictor \\
Mechanism & Deterministic scaling & Probabilistic association \\
Signal location & Mean structure & Covariance structure \\
Carryover impact & Modest (up to 7\%) &
  Design-dependent, up to 31\% \\
Appropriate when &
  \parbox[t]{3cm}{Biomarker\\determines dose\\or PK} &
  \parbox[t]{3cm}{Biomarker\\indexes a latent\\subtype} \\
\bottomrule
\end{tabular}
\end{table}

The choice of DGP architecture should be guided by the hypothesized
biological mechanism, as summarized above. When the mechanism is
ambiguous, as it often is in early biomarker development, we
recommend that investigators run the simulation under both
architectures and report the range of power estimates, treating
covariance moderation as the conservative default because its
larger carryover sensitivity keeps power estimates from being
optimistic. A dual-channel sensitivity sweep across mixing weights
is available in the companion paper\textsuperscript{\citeproc{ref-pmsimstats-paper11}{6}}.
Designing the trial to minimize carryover, through adequate washout
periods, also reduces sensitivity to the architecture choice. More
generally, simulation studies evaluating biomarker-moderated
treatment effects should state explicitly the following.

\begin{itemize}\setlength{\itemsep}{2pt}\setlength{\parskip}{0pt}
\item Whether the interaction is generated through mean moderation,
covariance moderation, or both channels simultaneously.
\item The biological rationale for the chosen mechanism.
\item Whether any carryover model operates on the interaction signal
consistently with that choice.
\item When the mechanism is uncertain, sensitivity analyses under the
alternative architecture.
\end{itemize}

These requirements specialize the Data-generating
element of the ADEMP framework\textsuperscript{\citeproc{ref-morris2019simulation}{8}} to the class
of architecture-sensitive biomarker simulations.

\subsection{4.3 Possible mitigation strategies}

Under covariance moderation, the power loss has two sources. The
drug-exposure variable's variance shrinks (recoverable), and the
biomarker-response correlation genuinely fades off drug (lost
information no model can recover). Because off-drug points carry
less signal, methods that down-weight or drop them may recover part
of the loss. Candidates include restricting the analysis to on-drug points, weighting
observations by expected residual drug effect (via the
\texttt{weights} argument in \texttt{nlme::lme}), a within-subject
contrast regressing each patient's on-minus-off mean difference on
the biomarker, a two-stage random-slopes approach, or excluding the
most contaminated early off-drug points. None of these strategies
was tested by simulation here. A companion paper on
carryover-mitigation strategies takes that up. We regard a
systematic simulation comparison, analogous to the factorial
comparison of Section 3, as important future work, and expect the
on-drug-only restriction and the weighted analysis to be the most
tractable candidates for a first empirical evaluation.

\subsection{4.4 Relation to the wider literature}

Almost every trial-simulation framework uses mean moderation.
Simon\textsuperscript{\citeproc{ref-simon2010}{2}}, Freidlin and Korn\textsuperscript{\citeproc{ref-freidlin2014}{3}}, Renfro and Sargent\textsuperscript{\citeproc{ref-renfro2017}{17}}, and standard N-of-1 simulation studies all put the
biomarker effect in the mean. Even methodological work on N-of-1
carryover sharing authors with Hendrickson et al.\textsuperscript{\citeproc{ref-hendrickson2020}{4}}
uses mean moderation, which suggests Hendrickson's
differential-correlation choice was a one-off rather than a
convention. The practical worry is that most published power
estimates may be optimistic for biomarkers that really behave like
covariance moderation, because they ignore the extra
carryover-driven loss demonstrated here. The distinction matters
most exactly in the designs precision medicine favors:
treatment-switching designs (crossover, N-of-1 hybrid,
basket/platform trials) put patients through off-drug phases where
the covariance-moderation correlation decays, so the interaction
signal in later periods is weaker than in the first. Parallel-group
enrichment designs, with no off-drug phase, are immune, and give
similar power under either architecture. Adaptive enrichment trials
that use interim power calculations written in the mean-moderation
idiom are at particular risk. If the truth is closer to covariance
moderation and patients pass through washouts, the real power is
lower than the rule assumes, which can trigger premature stopping
or an over-narrow eligible population. This is a
data-generating-process question, separate from, and not answered
by, the analysis-model guidance in the PATH statement\textsuperscript{\citeproc{ref-kent2020path}{7}}. The architecture choice matters most in exactly
the within-subject designs precision medicine finds most
attractive, and least in the parallel-group designs where it is
often, wrongly, assumed to be harmless.

\subsection{4.5 Limitations}

Several limitations qualify these results. The magnitude of the
covariance-moderation loss in the focal Hybrid design (15.4\%),
and the largest value observed across the three designs, the
30.6\% loss in OL+BDC, both fall below the 40-60\% range suggested
by earlier exploratory runs in this project, so the precise figures
should be read as design- and calibration-specific rather than
universal. A confirmatory run at \(N = 140\) on the Hybrid design,
analogous to the \(N = 140\) OL+BDC check reported in an earlier
version of this analysis, has not yet been re-executed under the
current N-matching convention (see Reproducibility) and remains
future work. The uniformly
sub-nominal Type I error we observe reflects a known conservatism
of the \texttt{corCAR1} Wald test that is common-mode across
architectures but not itself resolved here. A cluster-robust remedy
is characterized in a companion calibration study\textsuperscript{\citeproc{ref-pmsimstats-paper10}{11}}. Nominal coverage was inferred from the
correspondence between empirical and model-based standard errors
rather than verified by a dedicated coverage simulation. The
mitigation strategies of Section 4.3 are mechanistic hypotheses,
not simulation-tested recommendations. Finally, the combined
architecture that activates both channels at once is deferred
entirely to a companion paper\textsuperscript{\citeproc{ref-pmsimstats-paper11}{6}}, and our
results are matched to a single reference application
(prazosin/PTSD) and three within-subject designs. Generalization to
other biomarkers, designs, or effect sizes should be checked rather
than assumed.

\bigskip\noindent

\rule{\textwidth}{0.4pt}\bigskip

\section{5. Conclusions}

\begin{enumerate}
\item The choice between direct mean moderation and covariance
  moderation for generating biomarker-treatment interactions in
  clinical trial simulations has substantive consequences for power
  estimates under carryover.

\item Covariance moderation is intuitively appealing whenever a
  biomarker is better understood as predictive than causal, and
  Section 4.1 sets out the reasons this is often the more defensible
  specification. That appeal carries a measurable cost under
  carryover, and the cost is design-dependent because it is driven by
  a mechanism (Section 3.2) that is itself design-dependent: carryover
  erodes the biomarker-response correlation that carries the
  covariance-moderation signal only during uninterrupted off-drug
  time, so the loss scales with how long each design sustains an
  off-drug period. In the focal Hybrid design, carryover reduces
  power by 15.4\% under covariance moderation against 6.9\% under
  mean moderation, roughly double. The crossover and
  open-label-plus-blinded-discontinuation contrast designs bracket
  this from opposite sides, from 3.0\% (crossover, not distinguishable
  from zero, because within-subject AR(1) correlation already carries
  most of the signal) to 30.6\% (OL+BDC, whose sustained blinded
  off-drug phase gives the correlation-erosion mechanism the most
  uninterrupted time to act).

\item The choice should be guided by the hypothesized biological
  mechanism. Mean moderation is appropriate for biomarkers that directly determine
  drug pharmacokinetics or pharmacodynamics, covariance moderation
  for biomarkers that statistically predict treatment response
  through latent subtype membership.

\item For the prazosin-PTSD application, with blood pressure as the
  predictive biomarker, covariance moderation is the more
  appropriate model, and the resulting power sensitivity to
  carryover should inform trial-design decisions about off-drug
  period duration.

\item The existing clinical trial simulation literature uses mean
  moderation almost exclusively. The Hendrickson et al. (2020) MVN
  approach is, to our knowledge, unique in the N-of-1 context. This
  means published power estimates for biomarker-moderated crossover
  and N-of-1 designs may be systematically optimistic for
  biomarkers that operate through a correlation-based mechanism.

\item Simulation studies for biomarker-moderated treatment effects
  should explicitly report their DGP architecture and, when the
  biological mechanism is uncertain, consider sensitivity analyses
  under the alternative.
\end{enumerate}

\bigskip\noindent

\rule{\textwidth}{0.4pt}\bigskip

\section{Reproducibility}

\(N\) denotes the total number of participants across all
randomization paths, and Section 3.1 runs every design at the same
total, \(N = 70\), allocated across that design's paths, so the
cross-design comparison in Section 3 is a design comparison rather
than a sample-size comparison. An earlier version of this analysis
fixed the per-path count at 35 instead of the total, which gave the
four-path Hybrid design 140 participants against 70 for the
two-path designs and confounded any Hybrid advantage with twice the
sample. Section 3.1 has been re-run at matched totals, and only the
Hybrid figures changed. A confirmatory run at \(N = 140\) on the
Hybrid design, to check that the architecture-dependent ordering
survives the power saturation that comes with a larger sample, is
planned but has not yet been executed under the current
N-matching convention and is not reported here. A structured,
section-by-section summary of the argument, without the
derivations, tables, or citations, is provided as a supplementary
document (\texttt{bullets.Rmd}).

\section{References}\label{references}

\protect\phantomsection\label{refs}
\begin{CSLReferences}{0}{1}
\bibitem[\citeproctext]{ref-pmsimstats-paper08}
\CSLLeftMargin{1. }%
\CSLRightInline{{pmsimstats team}. Test-procedure choice for the biomarker-treatment interaction in aggregated {N}-of-1 trials: Dichotomization, classical {RM-ANOVA}, mixed-effects, and generalized estimating equations. Published online 2026.}

\bibitem[\citeproctext]{ref-simon2010}
\CSLLeftMargin{2. }%
\CSLRightInline{Simon R. Clinical trial designs for evaluating the medical utility of prognostic and predictive biomarkers in oncology. \emph{Personalized Medicine}. 2010;7(1):33-47.}

\bibitem[\citeproctext]{ref-freidlin2014}
\CSLLeftMargin{3. }%
\CSLRightInline{Freidlin B, Korn EL. Biomarker enrichment strategies: Matching trial design to biomarker credentials. \emph{Nature Reviews Clinical Oncology}. 2014;11(2):81-90.}

\bibitem[\citeproctext]{ref-hendrickson2020}
\CSLLeftMargin{4. }%
\CSLRightInline{Hendrickson RC, Thomas RG, Schork NJ, Raskind MA. Optimizing aggregated {N}-of-1 trial designs for predictive biomarker validation: Statistical methods and practical considerations. \emph{Frontiers in Digital Health}. 2020;2:13. doi:\href{https://doi.org/10.3389/fdgth.2020.00013}{10.3389/fdgth.2020.00013}}

\bibitem[\citeproctext]{ref-pmsimstats-paper03}
\CSLLeftMargin{5. }%
\CSLRightInline{{pmsimstats team}. Latent-class and mixture-model formulations for biomarker-treatment interaction in {N}-of-1 trials. Published online 2026.}

\bibitem[\citeproctext]{ref-pmsimstats-paper11}
\CSLLeftMargin{6. }%
\CSLRightInline{{pmsimstats team}. A combined data-generating architecture for biomarker-treatment interactions: Channel super-additivity and mixing-weight sensitivity in aggregated {N}-of-1 trials. Published online 2026.}

\bibitem[\citeproctext]{ref-kent2020path}
\CSLLeftMargin{7. }%
\CSLRightInline{{Kent DM, Paulus JK, Klaveren D van, et al.} The predictive approaches to treatment effect heterogeneity ({PATH}) statement. \emph{Annals of Internal Medicine}. 2020;172(1):35-45.}

\bibitem[\citeproctext]{ref-morris2019simulation}
\CSLLeftMargin{8. }%
\CSLRightInline{Morris TP, White IR, Crowther MJ. Using simulation studies to evaluate statistical methods. \emph{Statistics in Medicine}. 2019;38(11):2074-2102. doi:\href{https://doi.org/10.1002/sim.8086}{10.1002/sim.8086}}

\bibitem[\citeproctext]{ref-pmsimstats-paper02}
\CSLLeftMargin{9. }%
\CSLRightInline{{pmsimstats team}. Carryover representation and inference for biomarker-treatment interaction tests in aggregated {N}-of-1 trials under covariance moderation: A factorial simulation comparison of nine analysis specifications. Published online 2026.}

\bibitem[\citeproctext]{ref-pmsimstats-paper06}
\CSLLeftMargin{10. }%
\CSLRightInline{{pmsimstats team}. Three-component decomposition of treatment response in aggregated {N}-of-1 trials: Pharmacological, expectation-driven, and natural-history components. Published online 2026.}

\bibitem[\citeproctext]{ref-pmsimstats-paper10}
\CSLLeftMargin{11. }%
\CSLRightInline{{pmsimstats team}. Type {I} calibration of fixed-effect interaction tests in linear mixed models: A staged diagnosis of working-covariance misspecification and a cluster-robust remedy. Published online 2026.}

\bibitem[\citeproctext]{ref-senn2016mastering}
\CSLLeftMargin{12. }%
\CSLRightInline{Senn S. Mastering variation: Variance components and personalised medicine. \emph{Statistics in Medicine}. 2016;35(7):966-977. doi:\href{https://doi.org/10.1002/sim.6739}{10.1002/sim.6739}}

\bibitem[\citeproctext]{ref-raskind2013}
\CSLLeftMargin{13. }%
\CSLRightInline{{Raskind MA, Peterson K, Williams T, et al.} A trial of prazosin for combat trauma {PTSD} with nightmares in active-duty soldiers returned from iraq and afghanistan. \emph{American Journal of Psychiatry}. 2013;170(9):1003-1010.}

\bibitem[\citeproctext]{ref-bauercurran2003}
\CSLLeftMargin{14. }%
\CSLRightInline{Bauer DJ, Curran PJ. Distributional assumptions of growth mixture models: Implications for overextraction of latent trajectory classes. \emph{Psychological Methods}. 2003;8(3):338-363.}

\bibitem[\citeproctext]{ref-lubkemuthen2005}
\CSLLeftMargin{15. }%
\CSLRightInline{Lubke GH, Muthén B. Investigating population heterogeneity with factor mixture models. \emph{Psychological Methods}. 2005;10(1):21-39.}

\bibitem[\citeproctext]{ref-nylund2007}
\CSLLeftMargin{16. }%
\CSLRightInline{Nylund KL, Asparouhov T, Muthén BO. Deciding on the number of classes in latent class analysis and growth mixture modeling: A {M}onte {C}arlo simulation study. \emph{Structural Equation Modeling}. 2007;14(4):535-569.}

\bibitem[\citeproctext]{ref-renfro2017}
\CSLLeftMargin{17. }%
\CSLRightInline{Renfro LA, Sargent DJ. Statistical controversies in clinical research: Basket trials, umbrella trials, and other master protocols. \emph{Annals of Oncology}. 2017;28(1):34-43.}

\end{CSLReferences}

\end{document}
